\documentclass[10pt,a4paper]{amsart}
\usepackage[margin=19mm]{geometry}
\usepackage{amsmath,amssymb,mathtools}
\usepackage[scr=rsfs,cal=euler]{mathalfa}
\usepackage{xfrac}
\usepackage[unicode,hidelinks,bookmarks=false]{hyperref}
\newcommand{\ie}{i.e.\ }
\newcommand{\cf}{cf.\ }
\newcommand{\resp}{respectively\ }
\newcommand{\Subsection}{\subsection}
\providecommand{\nobreakdash}{\nobreak\hbox{-}}
\setlength{\emergencystretch}{2em}
\multlinegap0pt
\usepackage{amssymb}
\usepackage{bbm}
\newtheorem{proposition}{Proposition}
\newtheorem{cor}{Corollary}
\newtheorem{theorem}{Theorem}
\usepackage{cleveref}
\crefname{proposition}{Proposition}{Propositions}
\crefname{cor}{Corollary}{Corollarys}
\crefname{theorem}{Theorem}{Theorems}
\newcommand{\mA}{\mathcal{A}}
\newcommand{\mB}{\mathcal{B}}
\newcommand{\mC}{\mathcal{C}}
\newcommand{\mD}{\mathcal{D}}
\newcommand{\mcal}{\mathcal}
\title{On stability of distance under some tensor products and some calculations}
\author{Sumit Kumar}
\date{Source: arXiv:2601.05154v1 (CC BY 4.0)}
\begin{document}
\maketitle

\noindent\emph{Source and licence.} On stability of distance under some tensor products and some calculations. Version arXiv:2601.05154v1 (CC BY 4.0), distributed by arXiv under the Creative Commons Attribution 4.0 International licence, http://creativecommons.org/licenses/by/4.0/, as recorded in the arXiv licence metadata for this exact version and shown on the abstract page https://arxiv.org/abs/2601.05154v1. Full licence notice: \texttt{licenses/arxiv-2601.05154-CC-BY-4.0.txt}.\par\smallskip

\noindent\textit{Editorial scope.} Counted unit: the article's theorem Projective-Ch (source0.tex lines 634--644). The article's complete original proof of it is delivered with it (source0.tex lines 646--670, one \texttt{proof} environment of 25 lines). No mathematics is written, repaired or re-derived: the statement and the proof are the article's own lines, in the article's own notation, and no retained line is substituted or re-flowed.  Retained uncounted as the source-local context the proof needs: lines 541--546; lines 623--632.  Nothing was omitted from any retained span, so the package carries no omission record.  No retained line contains a citation, so the wrapper carries no bibliography and no bibliography entry.  No retained line contains a figure environment, an image-inclusion command or drawing code, so no image is delivered and the package declares no asset dependency.  Editorial formatting only, in addition: an AMS \texttt{amsart} preamble that restates the article's own declarations and macros used by the retained lines, namely the environments \texttt{proposition}, \texttt{cor}, \texttt{theorem} on separate counters, together with the standard packages those lines need.  The retained environments print in this document as Proposition 1, Corollary 1, Theorem 1 in order of appearance, whereas the article prints the same items with the numbering of its own full document.  The complete native source of the article as distributed in the arXiv e-print for this exact version is bundled unchanged as \texttt{sources/192324/source0.tex}.
\begin{proposition}{\label{unit-ball}}
Let $\mcal{X}$ and $\mcal{Y}$ be Banach spaces.  Then the closed unit
ball $B_{1}(\mcal{X} \otimes^{\gamma} \mcal{Y})$ of $\mcal{X}
\otimes^{\gamma} \mcal{Y}$ is the closed convex hull of the set
$B_{1}(\mcal{X})\otimes B_{1}(\mcal{Y})$.
\end{proposition}

\begin{cor}\label{D-unital}
 Let $\mC $ and $\mD $ be $C^*$-algebras and, $\mA$ and $\mB$ be
 $C^*$-subalgebras of $\mC$. If $\mD$ is unital, then
\[
d_{KK}(\mA, \mB)\leq d_{KK}(\mA \otimes^{\gamma} \mD, \mB \otimes^{\gamma} \mD)
\]
 and \[
d_{0}(\mA, \mB)\leq d_{0}(\mA \otimes^{\gamma} \mD, \mB \otimes^{\gamma} \mD).
\]
\end{cor}

\begin{theorem}\label{Projective-Ch}
Let $\mC$ and $\mD$ be $C^*$-algebras and, $\mA$ and $\mB$ be
$C^*$-subalgebras of $\mC$. If $\beta$ is a positive scalar such that
$\mA \subseteq_{\beta}\mB$, then $\mA \otimes^{\gamma} \mD
\subseteq_{\beta}\mB \otimes^{\gamma} \mD$.  In particular,
$$d_0(\mA \otimes^{\gamma} \mD, \mB \otimes^{\gamma} \mD) \leq
d_0(\mA, \mB).$$

Moreover, if $\mD$ is unital, then
$$d_{0}(\mA \otimes^{\gamma} \mD, \mB \otimes^{\gamma} \mD) = d_{0}(\mA, \mB).$$ 
\end{theorem}

\begin{proof}
As $B_{1}(\mA \otimes \mD)$ is dense in $B_{1}(\mA \otimes^{\gamma}
\mD)$, it is enough to show that $\mA \otimes \mD \subseteq_{\beta}\mB
\otimes \mD$.

Let $x\in B_{1}(\mA \otimes \mD)$. Then, by \Cref{unit-ball} $x=
\sum_{i=1}^{n}\alpha_{i} (a_{i}\otimes d_{i})$, where $a_{i}\in
B_{1}(\mA)$, $d_{i}\in B_{1}(\mD)$, $0\leq \alpha_{i}\leq 1$ and
$\sum_{i=1}^{n}\alpha_{i}= 1$. As $\mA \subseteq_{\beta} \mB$, for
each $a_{i}$, there exists a $b_{i}$ in $\mB$ such that
\(
\|a_{i}-b_{i}\| \leq  \beta.
\) 
Thus, $y := \sum_{i=1}^{n}\alpha_{i} (b_{i}\otimes d_{i})\in \mB\otimes \mD$ and 
\[
\|x-y\|_{\gamma} = 
\|\sum_{i=1}^{n}\alpha_{i} (a_{i}-b_{i})\otimes d_{i}\|_{\gamma} \leq  \sum_{i=1}^{n}\alpha_{i} \|a_{i}-b_{i}\|\|d_{i}\|  \leq
  \beta.
\]
Hence, $\mA \otimes \mD \subseteq_{\beta}\mB
\otimes \mD$.

When $\mcal{D}$ is unital, then the final equality follows from
\Cref{D-unital}.
\end{proof}

\end{document}
